% Build: module load texlive/20240312-GCC-13.3.0 && pdflatex top.tex (twice)
% Facts, all under /nfs/roberts/scratch/pi_nrs36/jar335/rebuild/explore/2026-10-29/:
%   top_research/PLAN.md (the approved plan), top_research/REPORT.md (sections 1, 2, 5),
%   top_research/REVIEW_PLAN_codex.md (BSYZ ratios, giving and deduction figures, the 2019 Forbes list),
%   top_research/pwbm/pwbm.txt (section 2.1), top_build/hidden/RESULTS.md (Part A results).
% SCF training rows: research/rebuild/FEATURES.md, Data (138,902 rows less the 2% held back).
% Part B and C results: top_build/billionaires/RESULTS.md.

% Shared look, copied from ../2026-10_rebuild/rebuild_talk.tex (lines 1-75)

\documentclass[11pt,aspectratio=169]{beamer}

% ---------------------------------------------------------------------------
% Look: plain beamer, Yale colours
% ---------------------------------------------------------------------------
\usetheme{default}
\usefonttheme{professionalfonts}
\setbeamertemplate{navigation symbols}{}

\definecolor{yaleblue}{HTML}{00356B}
\definecolor{yalemid}{HTML}{286DC0}
\definecolor{yalelight}{HTML}{63AAFF}
\definecolor{yalegrey}{HTML}{4A4A4A}
\definecolor{yalelightgrey}{HTML}{DDDDDD}

\setbeamercolor{structure}{fg=yaleblue}
\setbeamercolor{normal text}{fg=black}
\setbeamercolor{frametitle}{fg=yaleblue}
\setbeamercolor{title}{fg=yaleblue}
\setbeamercolor{alerted text}{fg=yalemid}
\setbeamercolor{block title}{fg=white,bg=yaleblue}
\setbeamercolor{block body}{bg=yalelightgrey!40}

\setbeamerfont{frametitle}{series=\bfseries,size=\large}
\setbeamerfont{title}{series=\bfseries,size=\Large}

\setbeamertemplate{itemize items}[circle]
\setbeamertemplate{blocks}[rounded][shadow=false]

% A thin blue rule under each frame title
\setbeamertemplate{frametitle}{%
  \vspace*{0.6em}%
  \usebeamerfont{frametitle}\usebeamercolor[fg]{frametitle}\insertframetitle\par
  \vspace*{-0.2em}{\color{yaleblue}\rule{\textwidth}{0.6pt}}%
}

% Footer: short title on the left, slide number on the right
\setbeamertemplate{footline}{%
  \hspace*{1.2em}{\color{yalegrey}\scriptsize\insertshorttitle}%
  \hfill{\color{yalegrey}\scriptsize\insertframenumber}\hspace*{1.2em}%
  \vspace*{0.6em}%
}

% ---------------------------------------------------------------------------
% Packages
% ---------------------------------------------------------------------------
\usepackage[T1]{fontenc}
\usepackage{lmodern}
\usepackage{booktabs}
\usepackage{array}
\usepackage{graphicx}
\usepackage{amsmath}
\usepackage{tikz}
\usetikzlibrary{positioning,arrows.meta,calc}
\usepackage{pgfplots}
\pgfplotsset{compat=1.18}
\graphicspath{{figures/}}

% Diagram styles
\tikzset{
  tok/.style={draw=yaleblue, rounded corners=2pt, font=\small, inner sep=4pt,
              minimum height=0.65cm},
  blank/.style={tok, fill=yalelight!35, font=\small\bfseries},
  box/.style={draw=yaleblue, rounded corners=2pt, font=\scriptsize, align=center,
               text width=2.1cm, minimum height=1.1cm, inner sep=3pt},
  arr/.style={-{Stealth[length=5pt]}, yaleblue, thick},
  cF/.style={fill=yaleblue},
  cP/.style={fill=yalelight!60},
  cN/.style={fill=white},
}

\pgfplotsset{colormap={yalemap}{color=(white) color=(yalelight) color=(yaleblue)}}



\newcommand{\slidenote}[1]{\par\vfill{\scriptsize\color{yalegrey}#1}}

\title[The PUF's aggregate records and the billionaires]{The PUF's aggregate records and the billionaires\\in the Tax-Data rebuild}
\author{John Ricco}
\institute{The Budget Lab at Yale}
\date{October 2026}

\hyphenpenalty=10000
\exhyphenpenalty=10000

\begin{document}

\begin{frame}[plain,noframenumbering]
  \titlepage
\end{frame}

% ---------------------------------------------------------------------------
\begin{frame}{The PUF's aggregate records}
  \small
  The rebuild is a generative model of tax units, trained on the 2015 PUF and four surveys. Records
  drawn from it are reweighted to each year's published targets.
  \vspace{0.4em}
  \begin{itemize}
    \item From tax year 2009, SOI replaced the most extreme returns in the PUF with records holding their
          group means.
    \item In the 2015 PUF, the last one released, 1,233 returns sit in four such records by AGI group: the
          top 400 by AGI, plus any return with one of the largest amounts of some item.
  \end{itemize}
  \vspace{0.3em}
  \setlength{\tabcolsep}{3pt}
  \begin{columns}[T]
    \begin{column}{0.47\textwidth}
      \centering\scriptsize
      \begin{tabular}{lrrr}
        \toprule
        AGI group & Returns & Mean AGI & Total AGI \\
        \midrule
        Negative          & 140   & $-$\$96.4m & $-$\$13.5bn \\
        \$0--10 million   & 233   & \$4.0m     & \$0.9bn \\
        \$10--100 million & 394   & \$43.8m    & \$17.2bn \\
        Over \$100 million & 465  & \$291.8m   & \$135.7bn \\
        \midrule
        All four          & 1,233 &            & \$140.3bn \\
        \bottomrule
      \end{tabular}
    \end{column}
    \begin{column}{0.51\textwidth}
      \centering\scriptsize
      \begin{tabular}{lrr}
        \toprule
        AGI above \$59.4 million & Returns & AGI \\
        \midrule
        SOI Table 4.3, top 0.001\% & 1,412 & \$214.6bn \\
        Retained PUF records       & 840   & \$71.1bn \\
        \bottomrule
      \end{tabular}
      \vspace{0.4em}

      \raggedright
      The largest retained AGI, \$136.2 million, is below the group's mean, \$152.0 million. No
      reweighting of retained records can reach both the count and the AGI.
    \end{column}
  \end{columns}
  \slidenote{2015 dollars. Table 4.3: SOI's table of returns at the top percentiles of AGI.}
\end{frame}

% ---------------------------------------------------------------------------
\begin{frame}{The billionaires, and handling both problems after generation}
  \footnotesize
  \begin{itemize}
    \item \textbf{No source has the billionaires.} The SCF removes the Forbes 400 by name before it draws
          its sample. The PUF records no wealth, and it hides the richest returns in the aggregate records.
    \vspace{0.4em}
    \item \textbf{Adding billionaire records to the SCF before training would not work}, for three reasons.
    \begin{enumerate}
      \item A few hundred rows among about 136,000 SCF training rows would barely move the network, so
            generation would almost never draw anything like them.
      \item The list gives wealth but no tax items, so the network would have nothing to learn their
            income from.
      \item The network predicts a size band for each amount, and the dollar figure is a real training
            record's value from that band. Nothing above the training data's largest values can be drawn.
    \end{enumerate}
    \vspace{0.4em}
    \item \textbf{So both problems are handled after generation and before the reweighting.} Pretraining
          does not change, and there is no retrain.
  \end{itemize}
  \slidenote{SCF training rows: four waves, each household stored five times by the SCF (``implicates'').
             The exclusion: Bricker, Henriques, Krimmel and Sabelhaus (2016), p.~280.}
\end{frame}

% ---------------------------------------------------------------------------
\begin{frame}{The hidden returns rebuilt from generated records (Part A)}
  \footnotesize
  \begin{enumerate}\setlength{\itemsep}{1pt}
    \item \textbf{Candidates.} Draw 510,000 extra 2015 PUF records from the networks, oversampling high AGI
          and large losses. Each keeps its population weight (one over its chance of being drawn).
    \item \textbf{Selection by SOI's rule} (Bryant, Czajka, Ivsin and Nunns 2014): the top 400 by AGI, the 30
          highest of each income item, the 30 lowest of each loss, the 10 highest or lowest of most other fields,
          ranked in population weight. SOI has not published the rule in full.
    \item \textbf{Four AGI groups}, as SOI's. Starting weights are each group's count from the raw records
          (140.31, 233.47, 394, 465), split equally over its records.
    \item \textbf{Counts.} Move the weights as little as possible so that the number of returns with a
          positive and with a negative amount of each field matches SOI's booklet table.
    \item \textbf{Dollars.} Scale each item so that each group's net amount equals its raw record and the
          positive and negative totals equal the booklet's. Totals are recomputed from their parts.
    \item \textbf{Wealth} drawn again from the network given the new incomes, below \$1 billion.
    \item \textbf{Growth} to the target year like the rest of the file. Table 4.3's top 0.001\% is a target again.
  \end{enumerate}
  \slidenote{Steps 1--6 are in 2015 dollars. The method follows Ricco (2020), Penn Wharton Budget Model, section 2.1,
             which started from 134 retained 2012 PUF records.}
\end{frame}

% ---------------------------------------------------------------------------
\begin{frame}{The billionaires built from the Forbes list (Part B)}
  \scriptsize
  \setbeamerfont{block title}{size=\footnotesize}
  \setlength{\itemsep}{0pt}
  \begin{columns}[T]
    \begin{column}{0.53\textwidth}
      \begin{itemize}\setlength{\itemsep}{1pt}
        \item \textbf{The list.} Forbes World's Billionaires, March 2019: 607 US billionaires,
              \$3.1 trillion. Spouses become one joint return; a family fortune stays one record.
        \item \textbf{The SCF's own units at \$1 billion or more are dropped}; SCF wealth targets are
              rebuilt below \$1 billion. A generated record at \$1 billion or more has its balance sheet
              drawn again.
        \item \textbf{Wealth.} Listed public holdings are public stock; the rest is a private business
              (real estate for real-estate fortunes). No home, no debt. Basis zero if self-made, equal to
              value if inherited.
        \item \textbf{Each rank group's income} from Balkir, Saez, Yagan and Zucman (2025): reported
              income is 1.75\% of wealth for the top 100 and 3.46\% for ranks 101 and below (their
              next-300 figure), split by kind as in their table.
        \item \textbf{Charity}: gifts of 0.7\% and 0.6\% of wealth; the calculator applies the limits.
        \item Age, sex and marital status from Forbes. \textbf{Weight 1, held fixed.}
      \end{itemize}
    \end{column}
    \begin{column}{0.44\textwidth}
      \begin{block}{Each person's share of the group's income}
        In proportion to the matching asset:
        \begin{itemize}\setlength{\itemsep}{0pt}
          \item dividends and capital gains: public stock;
          \item business profits, and separately losses: private holdings;
          \item interest, other income, wages and pensions: total wealth.
        \end{itemize}
        \vspace{0.1em}
        \textit{Example, with illustrative figures.} A group holds \$400 billion of public stock, of which Zuckerberg holds \$60
        billion, or 15\%. If the group's dividends are \$5 billion, his are \$0.75 billion. A founder of a
        wholly private company gets no dividends but a share of business profits.
      \end{block}
      Each kind is then split into tax items (for example, short- and long-term gains) as in Part A's
      rebuilt over-\$100-million group.
    \end{column}
  \end{columns}
\end{frame}

% ---------------------------------------------------------------------------
\begin{frame}{One reweighting of the whole file (Part C)}
  \footnotesize
  \begin{enumerate}
    \item \textbf{Billionaires held at weight 1.} Their contribution is subtracted from each target, and
          only the other weights are solved for.
    \vspace{0.3em}
    \item \textbf{A check before solving.} What remains of each target after the billionaires must be
          reachable with nonnegative weights on the other records. Any target the billionaires alone
          exceed is reported.
    \vspace{0.3em}
    \item \textbf{The reweighting} uses the usual targets plus Table 4.3's top 0.001\%. Every other record
          can move, the rebuilt hidden returns included. A target is met if it is within one standard
          error (half a standard error for SCF targets).
  \end{enumerate}
  \vspace{0.6em}
  \textbf{Why nothing is counted twice}
  \begin{itemize}
    \item \textbf{Income.} SOI's totals already include billionaires' income. With the billionaires fixed,
          the reweighting has to lower other top records to hit SOI's top totals. Among them are the rebuilt
          hidden returns, which partly stand for the same people.
    \item \textbf{Wealth.} The SCF targets cover the population below \$1 billion, and the list supplies
          everything above it.
  \end{itemize}
\end{frame}

% ---------------------------------------------------------------------------
\begin{frame}{Data sources}
  \scriptsize
  \renewcommand{\arraystretch}{1.15}
  \begin{tabular}{>{\raggedright}p{3.5cm} >{\raggedright\arraybackslash}p{9.6cm}}
    \toprule
    Source & What it gives, and where it is used \\
    \midrule
    Balkir, Saez, Yagan and Zucman (2025), and the IRS study by the same authors (2024)
      & Forbes 400 matched to tax returns, means per tax unit over 2018--2020. ``Fiscal income'' is market
        income reported on returns, including tax-exempt interest and excluding prior-year net operating
        losses. Fiscal income over wealth: 1.7515\% for the top 100 (92 tax units), 3.4627\% for the next
        300 (276). Shares in gains, dividends, interest, business, other, wages and pensions: 62, 28, 11,
        $-$11, 6, 5\% and 60, 16, 8, 4, 6, 5\% (the underlying amounts are used). Gifts 0.7\% and 0.6\% of
        wealth. Deduction check: \$105.3m deducted of \$140.1m given per top-group unit (75\%); about 94\%
        for the next group. (Part B) \\
    SOI, 2015 PUF booklet, printed pp.~32--36
      & For all four aggregate records together: returns and dollars with a positive and with a negative
        amount, 171 fields. (Part A) \\
    Forbes World's Billionaires, March 2019
      & 607 US billionaires, \$3.1 trillion. (Part B) \\
    SOI Table 4.3
      & Returns, AGI and items at the top 0.001\%, 0.01\%, 0.1\% and 1\% by AGI. (Part A check; targets) \\
    Ricco (2020), ``Data Processing for PWBM's Tax Module,'' section 2.1
      & The method Part A follows: copy the extreme records, rescale weights to each group's count, match
        counts as closely as feasible, then scale dollars by AGI group. (Part A) \\
    \bottomrule
  \end{tabular}
\end{frame}

% ---------------------------------------------------------------------------
\begin{frame}{How other models handle the very top}
  \scriptsize
  \renewcommand{\arraystretch}{1.25}
  \begin{tabular}{>{\raggedright}p{3.0cm} >{\raggedright\arraybackslash}p{10.1cm}}
    \toprule
    \multicolumn{2}{l}{\textbf{The aggregate records}} \\
    PolicyEngine & Since March 2026, each record is rebuilt as 20--40 synthetic returns started from
                   retained returns in the same AGI group, calibrated to the record's totals. The
                   over-\$100-million group is built from a Forbes 400 snapshot. \\
    Tax Policy Center & Replaces them with several thousand extreme records from the 2001--2007 PUFs,
                        reweighted to the aggregate totals (from its documentation as quoted elsewhere;
                        not read directly). \\
    PSL taxdata and tmd & Drop them. \\
    CBO, JCT, Treasury & Use SOI's confidential sample, which keeps every high-income return. \\
    \midrule
    \multicolumn{2}{l}{\textbf{Rich lists}} \\
    Federal Reserve (DFA) & Adds the Forbes 400 to the SCF, with income from the top 0.1\%'s median
                            income-to-wealth ratio. The paper splits weight between the SCF and Forbes in
                            wealth bins; the public code adds the list at weight 1. \\
    CBO (2024) & Places Forbes families at the top of the SCF, with portfolios from other families in the
                 top 0.1\% of wealth. \\
    Tax Policy Center & Matches each Forbes 400 member to a model record of the same age and marital status,
                        with assets from a matched SCF record. \\
    \bottomrule
  \end{tabular}
  \slidenote{Sources and page references: REPORT.md, section 5 (29 October 2026).}
\end{frame}

% ---------------------------------------------------------------------------
\begin{frame}{Results for the hidden returns, 2015 (Part A)}
  \footnotesize
  \centering
  \begin{tabular}{lrrr}
    \toprule
    Top 0.001\% by AGI, 2015 & Generated file & With rebuilt returns & SOI Table 4.3 \\
    \midrule
    Returns                  & 1,419     & 1,420     & 1,412 \\
    AGI floor                & \$49.1m   & \$61.8m   & \$59.4m \\
    AGI                      & \$103.2bn & \$216.1bn & \$214.6bn \\
    Dividends                & \$6.7bn   & \$23.3bn  & \$23.4bn \\
    Capital gains            & \$59.7bn  & \$122.6bn & \$118.3bn \\
    Partnership and S corporation income & \$20.2bn & \$38.0bn & \$37.7bn \\
    \bottomrule
  \end{tabular}
  \raggedright
  \begin{itemize}
    \item 2,212 rebuilt records stand for the 1,233 returns. All 64 booklet counts the rebuild carries are
          met within the larger of 2 returns and 2\%. Each group's AGI equals its raw record's.
    \item The wider groups overshoot before reweighting. The generated file alone already has SOI's AGI
          for the top 0.1\% (\$1,032.9bn against \$1,033.5bn); with the rebuilt returns it is 15\% above.
          The reweighting has to lower other top records.
    \item Every rebuilt return has wealth below \$1 billion. In the over-\$100-million group, median net
          worth rises from \$140 million to \$255 million when drawn again given the new incomes.
  \end{itemize}
  \slidenote{2015 dollars, except net worth (2019 dollars). Generated filers at their 2015 population
             weights, before any reweighting; dependent filers excluded, as in Table 4.3.}
\end{frame}

% ---------------------------------------------------------------------------
\begin{frame}{Results: the billionaires}
  \footnotesize
  \centering
  \renewcommand{\arraystretch}{1.1}
  \begin{tabular}{>{\raggedright}p{4.6cm} >{\raggedright}p{4.6cm} >{\raggedright\arraybackslash}p{3.4cm}}
    \toprule
    Check & Compared with & Result \\
    \midrule
    Count and wealth of the list & Forbes: 607 people, \$3.1 trillion & 607 people, \$3,111.0 billion; 605 tax units after joining two couples \\
    AGI per unit, top 92 units & IRS study (BSYZ): \$345 million & \$329 million \\
    AGI per unit, next 276 units & IRS study: \$115 million & \$124 million \\
    Gifts per unit, top group & IRS study: \$140 million & \$133 million \\
    Gifts per unit, next group & IRS study: \$22.0 million & \$21.7 million \\
    SCF units dropped at \$1 billion or more & Forbes ranks 401 and below: 205 units, \$298 billion & 360 units, \$430 billion \\
    \bottomrule
  \end{tabular}
  \raggedright
  \begin{itemize}
    \item The calculator has not been run, so the charitable deduction and federal income tax are not yet checked against the IRS study.
  \end{itemize}
  \slidenote{Nominal 2019 dollars. The test used the networks from the 28--29 October retrain.}
\end{frame}

% ---------------------------------------------------------------------------
\begin{frame}{Results: the reweighting}
  \footnotesize
  \centering
  \renewcommand{\arraystretch}{1.1}
  \begin{tabular}{>{\raggedright}p{5.2cm} >{\raggedright}p{4.0cm} >{\raggedright\arraybackslash}p{3.4cm}}
    \toprule
    Check & Compared with & Result \\
    \midrule
    Targets met & Every 2019 target, 805 in all & All met, with and without the billionaires \\
    Top 0.001\% by AGI & SOI Table 4.3, 2019: \$214.5 billion & \$212.3 billion \\
    Top 0.001\% returns & SOI: 1,482 & 1,492 \\
    Billionaire weights & Exactly 1 & Exactly 1, all 605 \\
    Effective sample of free records & File without billionaires: 80.33\% & 80.32\% \\
    Part A's over-\$100-million rebuilt returns & Their starting values: 470 returns, \$151 billion AGI & 336 returns, \$65 billion AGI \\
    Wealth below \$1 billion & SCF target: \$96.65 trillion & \$96.83 trillion; the list adds \$3.11 trillion \\
    \bottomrule
  \end{tabular}
  \slidenote{Nominal 2019 dollars. The test used the networks from the 28--29 October retrain.}
\end{frame}

% ---------------------------------------------------------------------------
\begin{frame}{Open questions and limits}
  \footnotesize
  \begin{itemize}
    \item \textbf{Inside the hidden groups, the joint distribution is assumed.} Other candidate records
          could match the same SOI totals and give different answers for policies with thresholds in this
          range, such as a minimum tax on very high incomes.
    \item \textbf{Family fortunes are kept as one record}, though a family may file several returns.
    \item \textbf{Ratios below rank 400 are extrapolated.} Balkir, Saez, Yagan and Zucman cover only the top 400; the next-300 ratio
          is used for everyone below.
    \item \textbf{Income shared by holdings is an approximation.} A private C corporation can pay
          dividends, and a gain can come from an asset no longer held.
    \item \textbf{Replacing the SCF's households above \$1 billion with the list is an assumption.} The
          Federal Reserve estimated that about \$300 billion of Forbes 400 wealth was already represented
          in the 2019 SCF.
    \item \textbf{Wealth for the rebuilt returns is an extrapolation.} The network sees their incomes only
          as its top bands, the same bands a \$20 million income falls in.
    \item \textbf{Later years} start from the 2015 rebuild, grown forward, and need each year's Table 4.3
          and each year's Forbes list. Charity and some deductions are not grown from 2015 dollars.
  \end{itemize}
\end{frame}

\end{document}
